%% Chapter 4, Section 4.1: Metropolis Hastings Algorithm

\section{Metropolis Hastings Algorithm}
\label{sec:4.1}

The Metropolis Hastings algorithm provides a flexible framework to define a Markov kernel $p_{\mathrm{MH}}$
that can be sampled from and leaves the target invariant. The idea is to leverage a probabilistic
accept/reject mechanism to turn a user-chosen proposal Markov kernel $q$ into a Markov kernel
$p_{\mathrm{MH}}$ that satisfies detailed balance with respect to the target $f$. From the definition of Markov
kernel, for each $x \in E$, $q(x, z)$ will represent the probability (or probability density in the
continuous case) with which a move from $x$ to $z$ is proposed. Recall further that, for fixed $x$,
$q(x, \cdot)$ is a p.m.f.\ if $E$ is discrete and a p.d.f.\ if $E$ is continuous.

The method is summarized in Algorithm~\ref{alg:4.1}. Given the $n$-th sample $X^{(n)}$, we first draw a
proposed move $Z^* \sim q(X^{(n)}, \cdot)$. We accept the move, which means setting $X^{(n+1)} = Z^*$,
with probability $a(X^{(n)}, Z^*)$. If the proposed draw $Z^*$ is rejected, we set $X^{(n+1)} = X^{(n)}$. The
probability of accepting a proposed move from $x \in E$ to $z \in E$ is defined to be
\[
  a(x, z) := \min\!\left\{1,\, \frac{f(z)\, q(z, x)}{f(x)\, q(x, z)}\right\}.
\]
We will show in Theorem~\ref{thm:4.3} that this choice of acceptance probability turns the kernel $q$ into a
kernel $p_{\mathrm{MH}}$ that satisfies detailed balance with respect to $f$.

\begin{algorithm}[H]
\caption{Metropolis Hastings Algorithm}
\label{alg:4.1}
\begin{algorithmic}[1]
\Require Target $f$, initial distribution $\pi_0$, proposal Markov kernel $q(x,z)$, sample size $N$.
\State Initial draw: Sample $X^{(0)} \sim \pi_0$.
\State Subsequent draws: For $n = 0, \ldots, N-1$ do:
\State \textbf{Proposal step:} Sample $Z^* \sim q(X^{(n)}, \cdot)$.
\State \textbf{Accept/reject step:} Update
\[
  X^{(n+1)} =
  \begin{cases}
    Z^*     & \text{with probability } a(X^{(n)}, Z^*),\\
    X^{(n)} & \text{with probability } 1 - a(X^{(n)}, Z^*).
  \end{cases}
\]
\Ensure Sample $\{X^{(n)}\}_{n=1}^{N}$.
\end{algorithmic}
\end{algorithm}

As previously mentioned, given a test function $h : E \to \R$, we can use the output $\{X^{(n)}\}_{n=1}^{N}$
of the Metropolis Hastings algorithm to approximate
\[
  \calI_f[h] = \int_E h(x) f(x)\, dx
  \approx \frac{1}{N} \sum_{n=1}^{N} h(X^{(n)}) =: \calI_f^{\mathrm{mh}}[h].
\]

\begin{remark}
\label{rem:4.1}
\begin{enumerate}
\item \textit{To implement the Metropolis Hastings algorithm we need to be able to sample $q(x, \cdot)$ for
$x \in E$ and evaluate the acceptance probability $a(x, z)$ for $x, z \in E$. A priori, the only
necessary condition on the proposal is that}
\[
  E \subset \bigcup_{x \in E} \mathrm{support}\bigl(q(x, \cdot)\bigr).
\]
\textit{Importantly, the target $f$ appears in the acceptance probability only as a ratio, and therefore the algorithm can be implemented even if $f$ is only known up to a normalizing constant.}

\item \textit{If $q(x, z) = q(z, x)$ for all $x, z \in E$, then $a(x, z) = \min\!\left\{1, \frac{f(z)}{f(x)}\right\}$.
Moves to regions of higher target density are always accepted, while moves to regions of lower but non-zero
target density are accepted with positive probability in order to ensure exploration of the
state space $E$. If $q$ is not symmetric, moves for which $q(z, x) \geq q(x, z)$ are favored, as
they are more likely to be undone by chance.}

\item \textit{The accept/reject step can be implemented as follows:}
\begin{itemize}
  \item \textit{Draw $U^* \sim \Unif(0, 1)$ independently from $Z^*$.}
  \item \textit{Update}
\[
  X^{(n+1)} =
  \begin{cases}
    Z^*     & \text{if } U^* < a(X^{(n)}, Z^*),\\
    X^{(n)} & \text{if } a(X^{(n)}, Z^*) < U^*.
  \end{cases}
\]
\end{itemize}
\end{enumerate}
\end{remark}

\medskip
\noindent\textbf{Pros and Cons.}
The Metropolis Hastings algorithm can be implemented without knowing the
normalizing constant of the target, and is extremely flexible due to the freedom in the choice
of proposal kernel $q$. It has important advantages over all the algorithms covered in Chapters
2 and 3. In contrast to rejection sampling, no samples are discarded; in contrast to importance
sampling, MCMC samples are given equal weights, avoiding weight collapse; in contrast
to ABC, in many cases MCMC is provably accurate in the large $N$ limit. The Metropolis
Hastings algorithm also has some caveats. First, it can be very hard (often impossible) to
tell if the chain is close to stationarity. As opposed to rejection sampling and importance
sampling, MCMC samples are typically positively correlated, which, as established in Lemma
A.36, results in larger variance estimates than independent or negatively correlated samples.
Additionally, it is often challenging to choose an appropriate proposal kernel, and the ergodic
behavior of the algorithm depends on that choice. In particular, it is often necessary to tune
the proposal by trial and error (the Gibbs sampler, studied in Chapter~5, is an exception).
Finally, the Metropolis Hastings algorithm cannot be naturally parallelized.

The Metropolis Hastings algorithm implicitly defines a Markov kernel $p_{\mathrm{MH}}(x, z)$ which gives
the probability (or probability density) of finding the $(n+1)$-th sample at location $z \in E$ given
that the $n$-th sample was at $x \in E$.

\begin{lemma}
\label{lem:4.2}
\textit{The Metropolis Hastings Markov kernel is given by}
\[
  p_{\mathrm{MH}}(x, z) = q(x, z)\, a(x, z) + \delta_x(z)\, r(x),
\]
\textit{where}
\[
  r(x) =
  \begin{cases}
    \displaystyle\sum_{y \in E} q(x, y)\bigl(1 - a(x, y)\bigr) & \textit{if } E \textit{ is discrete,}\\[8pt]
    \displaystyle\int_E q(x, y)\bigl(1 - a(x, y)\bigr)\, dy    & \textit{if } E \textit{ is continuous,}
  \end{cases}
\]
\textit{and $\delta_x(z)$ denotes a Dirac mass at $x$. (Note: $p_{\mathrm{MH}}$ is not absolutely continuous with respect to
Lebesgue measure when $E$ is continuous, since the probability of moving from $x$ to $x$ is positive.)}
\end{lemma}

\begin{proof}
We only prove the discrete case. The chain moves to a new state $z$ if the state $z$ was
proposed and accepted, which happens with probability $q(x, z) a(x, z)$. This is the probability of
moving from $x$ to $z$ if $x \neq z$. A move from $x$ to $x$ may occur in two different ways:
\begin{enumerate}
\item Propose $x$ as new state and accept it, which happens with probability $q(x, x) a(x, x)$.
\item Propose any $z \in E$ and reject it, which happens with probability
\[
  r(x) = \sum_{z \in E} q(x, z)\bigl(1 - a(x, z)\bigr).
\]
\end{enumerate}
Putting everything together
\[
  p_{\mathrm{MH}}(x, z) = q(x, z)\, a(x, z) + \delta_x(z)\, r(x). \qedhere
\]
\end{proof}

As part of the proof of the previous lemma, we have shown that if $x \neq z$, then
\[
  p_{\mathrm{MH}}(x, z) = q(x, z)\, a(x, z).
\]
A consequence of this identity is the detailed balance of $p_{\mathrm{MH}}$ with respect to $f$.

\begin{theorem}[Detailed Balance of Metropolis Hastings]
\label{thm:4.3}
\textit{The Metropolis Hastings kernel $p_{\mathrm{MH}}$ satisfies detailed balance with respect to $f$.}
\end{theorem}

\begin{proof}
We need to show that
\[
  f(x)\, p_{\mathrm{MH}}(x, z) = f(z)\, p_{\mathrm{MH}}(z, x), \quad \forall\, x, z \in E.
\]
For $x = z$ the equality is trivial. If $x \neq z$
\begin{align*}
  f(x)\, p_{\mathrm{MH}}(x, z)
  &= f(x)\, q(x, z)\, a(x, z) \\
  &= \min\!\left\{f(x)\, q(x, z),\, f(z)\, q(z, x)\right\}.
\end{align*}
The right-hand side is symmetric in $x$ and $z$, and therefore the result follows.
\end{proof}

Note that since detailed balance implies general balance, Theorem~\ref{thm:4.3} implies that $f$ is an
invariant distribution for $p_{\mathrm{MH}}$.
